\documentclass[11pt]{article}
\usepackage[margin=1in]{geometry}
\usepackage{lmodern}
\usepackage{amsmath}
\usepackage{graphicx}
\usepackage{hyperref}
\usepackage[numbers]{natbib}

\hypersetup{colorlinks=true, linkcolor=blue, urlcolor=blue, citecolor=blue}
\title{ARCMiS: Agentic Repository-Level Code Migration\\for Small Language Models}
\author{}
\date{September 2026}

\begin{document}
\maketitle

\begin{abstract}
We present ARCMiS (Agentic Repository-level Code Migration for SLMs), a
harness that migrates whole source codebases from one programming language
to another using small language models (SLMs, 7B--27B parameters) served
locally. The core method is a multi-agent system (MAS): a manager loop over
a ten-phase state machine, delegating to ten specialist roles through parsed
decision lines. Static agents meet dynamic selection---a task graph with
dependency-gated fan-out, team collectives for per-batch translation, a
deterministic tool guard on every call, per-role-class turn budgets, a
stagnation breaker, and a fleet model ladder. All state lives on a
character-capped shared blackboard. Long-context overflow is absorbed by
snapcompact, which renders discarded history to PNG bitmap frames instead
of calling an LLM summarizer. Scoring is toolchain-only: the target build
and test commands must pass in the produced workspace, and agent
self-reports never gate. In a 180-cell survey across three methods and
three SLMs, only 5 of 180 runs produced a toolchain-verified migration, and
all five came from the simplest method on the smallest project. A
Meta-Harness-style autooptimise loop diagnosed the binding failure of the
MAS baseline---reasoning-block truncation at the 4096-token output budget
stalling every run at the Contract phase---and a one-line budget fix
(max\_output\_tokens 4096 $\rightarrow$ 8192) let the first translated and
toolchain-verified project land. We report the architecture, the honest
negative results, and the failure taxonomy that disciplined the search.
\end{abstract}

\section{Introduction}
\label{sec:intro}

Repository-level code migration---translating an entire codebase from one
language to another while keeping its behavior---is one of the hardest open
tasks for coding agents. It needs three things at once: reading and
understanding code that does not fit in one context window, planning a
dependency-respecting translation order, and producing code that passes
real toolchains. Large frontier models can attempt this with ad-hoc
prompting; small language models cannot, and small models are the only
affordable option for sustained local use: they run on one consumer GPU,
cost nothing per token, and keep proprietary source on the machine.

ARCMiS targets exactly this regime. It migrates source codebases with
7B--27B models served locally through ollama, on a single RTX~3090. The
design constraint that shapes the whole system is \emph{honesty of
scoring}: an agent's claim of success counts for nothing. A run succeeds
only when the target toolchain compiles the produced workspace and its test
command exits zero. LLM-graded proxy scores, when recorded at all, never
gate.

This paper makes four contributions:
\begin{itemize}
\item A MAS methodology for repository-level migration with SLMs: ten
  phases, ten specialist roles, a shared blackboard, and a deterministic
  guard (Section~\ref{sec:method}).
\item Supporting infrastructure purpose-built for small local models:
  character-capped shared state, bitmap-frame context compaction, and
  hashline-based file invalidation (Section~\ref{sec:infra}).
\item Honest measurements: three method surveys (180 and 105 cells) and a
  Meta-Harness-style optimization loop with a diagnosed, mechanism-level
  fix (Sections~\ref{sec:results} and~\ref{sec:frontier}).
\item A failure taxonomy for SLM agents at this scale: reasoning-block
  truncation, provider 500s on continuation shapes, and stagnation
  (Section~\ref{sec:discussion}).
\end{itemize}

\section{Related Work}
\label{sec:related}

\paragraph{Reasoning-and-acting agents.}
ReAct~\cite{yao2023react} interleaves chain-of-thought with tool actions
and is the base loop for our monolith baseline. Reflexion~\cite{shinn2023reflexion}
adds verbal self-reflection across retries; Self-Refine~\cite{madaan2023selfrefine}
iterates a generate--feedback--revise loop. ARCMiS borrows the iterate-on-failure
shape but moves the feedback to deterministic judges and toolchain evidence,
because 27B models self-assess poorly.

\paragraph{Multi-agent frameworks.}
MetaGPT~\cite{hong2023metagpt} assigns a software company's roles (product
manager, architect, engineer) to LLM agents over standardized artifacts.
AutoGen~\cite{wu2023autogen} builds applications from conversing agents with
human-in-the-loop checkpoints. CAMEL~\cite{li2023camel} studies role-playing
communicative agents. These frameworks assume strong models; ARCMiS keeps
the role decomposition but replaces free-form conversation with a parsed
decision protocol, a static registry, and a state machine, because
conversation overhead is exactly what small turn budgets cannot afford.

\paragraph{Software-engineering agents.}
SWE-agent~\cite{yang2024sweagent} shows that a well-designed agent--computer
interface matters as much as the model; our hashline file tools and guard
are the same principle at small scale. Agentless~\cite{xia2024agentless}
demonstrates that a simple pipeline with localizer--repairer--validator
stages beats complex agents on some tasks; our survey results agree in the
small-project regime and our MAS adds dynamic scheduling for the parts that
do not fit a fixed pipeline.

\paragraph{Tool use.}
Toolformer~\cite{schick2023toolformer} teaches models to call APIs
self-supervised; Gorilla~\cite{patil2023gorilla} fine-tunes for API
invocation. ARCMiS uses neither: the 17-tool surface is fixed, and the
guard, not the model, enforces the calling discipline.

\paragraph{Code translation.}
TransCoder~\cite{lachaux2020transcoder} does unsupervised transpilation with
neural machine translation; AlphaTrans~\cite{arya2024alphatrans} is a
neuro-symbolic repository-level translator with behavioral checks.
CodaMosa~\cite{lemieux2023codamosa} escapes coverage plateaus by mixing LLM
test generation with search-based testing. Our problem sets reuse the
ReCodeAgent dataset's project pairs, and our validator--tester loop applies
the test-generation lesson with a per-module cap.

\paragraph{Small-model reasoning.}
DeepSeek-R1~\cite{deepseek2025r1} and Qwen3~\cite{qwen3_2025} make deliberate
reasoning (\texttt{think} blocks) standard on open models at our scale. This
is a double-edged capability: long think blocks consume output-token budgets
and directly caused the truncation failure we diagnose in
Section~\ref{sec:frontier}. Speculative decoding~\cite{leviathan2023speculative}
(and its MTP/SFT variants in the models we use) makes local inference fast
enough for long agent runs, which is why the 100-agent-hour budgets of this
paper fit on one GPU.

\paragraph{Harness optimization.}
Meta-Harness~\cite{metaharness2026} optimizes an executable harness through
an agentic outer loop where the proposer reads raw execution traces from the
filesystem. ARCMiS adopts its experiment contract wholesale
(Section~\ref{sec:frontier}); the ablation in that paper---raw traces beat
summaries for the proposer---is what motivated our immutable, greppable
experiment directories.

\section{The ARCMiS Methodology}
\label{sec:method}

\subsection{Overview}

A run is one migration attempt against one problem set: a source codebase in
language $L_s$ and a target language $L_t$, with a test command defined in
the config. The harness (a Rust binary built on rig~0.42~\cite{yao2023react})
spawns a MAS whose members share one blackboard directory. The manager
never writes product code; specialists never re-plan. Table~\ref{tab:roles}
shows the role taxonomy.

\begin{table}[ht]
\centering
\begin{tabular}{llp{7.2cm}}
\hline
Role & Class & Responsibility \\
\hline
Manager & control & Owns plan and task list; emits \texttt{DECISION} lines; no tools. \\
Analyst & worker & Reads the codebase; writes the source map and graphs. \\
Architect & worker & Writes the migration contract (\texttt{analysis/brief.md}). \\
Planner & worker & Produces dependency-ordered translation batches. \\
Translator & worker & Writes target-language modules for one batch. \\
Validator & judge & Judges one translated batch; re-runs the suite; read-only. \\
Tester & worker & Adds tests; runs the test command. \\
Failure-analyst & judge & Classifies one failure into a five-category taxonomy. \\
Critic & judge & End-of-run adversarial review; read-only. \\
Repairer & worker & Fixes one diagnosed failure. \\
Fleet-analyst & judge & Promotes or demotes roles along the model ladder. \\
\hline
\end{tabular}
\caption{The role taxonomy. Manager is tier 1; specialists are tier 0;
the translation collective is tier 2. Judge-class roles get small turn
budgets; worker-class roles get large ones.}
\label{tab:roles}
\end{table}

\subsection{Phase state machine}

The run walks ten phases: \textsf{Preflight}, \textsf{Discovery},
\textsf{Contract}, \textsf{Planning}, \textsf{Pilot}, \textsf{Migration},
\textsf{Integration}, \textsf{Hardening}, \textsf{FinalValidation},
\textsf{Done}. Forward transitions are earned: the manager's \texttt{done}
verb is accepted only when the phase's exit condition holds \emph{and} at
least one delegation completed in the phase (the \texttt{phase\_delegations}
counter). Regressions are legal and recorded: \textsf{Pilot} and
\textsf{Migration} may fall back to \textsf{Discovery} on a missing
dependency, \textsf{Integration} to \textsf{Migration}, and so on. Anything
else is a programming error, not a model decision. Each phase has a written
exit condition in the manager prompt (for example: \textsf{Contract}
requires \texttt{analysis/brief.md} to exist with a validator pass covering
it; \textsf{FinalValidation} requires the configured toolchain command to
exit zero on the target workspace).

\begin{table}[ht]
\centering
\begin{tabular}{llll}
\hline
\# & Phase & Exit evidence & Legal regressions \\
\hline
0 & Preflight & deterministic checks, no model call & --- \\
1 & Discovery & source map + brief cite every module & --- \\
2 & Contract & \texttt{brief.md} + validator pass & --- \\
3 & Planning & every module in one batch; order compiles & --- \\
4 & Pilot & one batch translated, validated, tested & Pilot$\to$Discovery \\
5 & Migration & every batch translated and validated & Migration$\to$Discovery \\
6 & Integration & full test suite green & Integration$\to$Migration \\
7 & Hardening & critic verdict is pass & Hardening$\to$Integration \\
8 & FinalValidation & toolchain rerun exits zero & FinalVal.$\to$Integration \\
9 & Done & --- & --- \\
\hline
\end{tabular}
\caption{The phase state machine. Forward moves are evidence-gated;
regressions are recorded, not hidden.}
\label{tab:phases}
\end{table}

\subsection{Decision protocol and dynamic scheduling}

The manager answers each round with plain text whose final lines carry
\texttt{DECISION:} verbs---\texttt{delegate}, \texttt{replan},
\texttt{escalate}, \texttt{done}. Parsing text instead of expecting tool
calls removes a whole failure class observed in the surveys: malformed
tool-call JSON that the ollama daemon rejects with HTTP 500 (13\% of all
220926 runs). The manager gets \emph{no tools at all}; the round prompt
carries the plan, notes, task list, ledger tail, and state, so the model
cannot spend turns re-reading files (an observed pathology: 40 consecutive
reads of \texttt{plan.md}).

Each \texttt{delegate} decision becomes a task node with optional
\texttt{after:} dependencies and an optional \texttt{team} clause. The task
graph executor runs the ready set (all dependencies done) with bounded
fan-out (default 2). Because ollama serves one model at a time, fan-out
interleaves turn streams rather than running parallel processes. Routing is
a three-stage chain: the manager's role hint, then a deterministic keyword
pass over the task text, then (as a last resort) the router model. After
each round the orchestrator re-scores the task list and the manager sees
merged state, so a failing task can spawn a failure-analyst $\rightarrow$
repairer sub-chain without any phase regression.

\paragraph{Team collectives.}
A \texttt{team=true} delegation runs the translation collective
translator $\rightarrow$ validator $\rightarrow$ tester, with
failure-analyst $\rightarrow$ repairer re-entry, as an inner loop under one
task id. The collective is tier 2, one level deep; the guard refuses deeper
nesting. It has its own turn budget (\texttt{collective\_turns}, default 40)
and stops early on a validator pass.

\subsection{Blackboard}

All mutable state lives in a per-run directory written by the blackboard
crate: \texttt{state.json} (phase, counters), \texttt{plan.md} and
\texttt{notes.md} (character-capped at 4000 and 8000), \texttt{tasks.json}
(task list mirror), \texttt{decisions.jsonl} and \texttt{failures.jsonl}
(ledgers), and a dependency graph. Character caps are the discipline that
keeps the shared state inside one read for a 16k-token context window; the
writer truncates with an elision marker rather than silently growing the
context of every later reader.

\subsection{Guard}

Every tool call from every agent passes one \texttt{on\_tool\_call} hook.
The deterministic tier checks: role allowlist (from static tool metadata),
path policy (\texttt{source/} is read-only; writes must stay inside the
workspace), deny patterns for bash (\texttt{rm -rf}, \texttt{git push},
package installs, network commands, \texttt{sudo}), and per-delegation
tool-call budgets. Rejections return \texttt{Skip(feedback)}: the model
sees the reason and self-corrects, and the breaker counts repeats. Calls the
deterministic tier marks ambiguous fall to an optional model-arbitration
tier (the ``Jev-as-guard'' slot), deny-by-default while only one model
exists. The guard is the single enforcement point; an earlier
wrapper-object design (a \texttt{GuardedTool} adapter) was deleted because
it was constructed but never on the call path.

\subsection{Budgets, breaker, fleet ladder}

Budgets are per role class: manager turns per round (default 20), worker
turns per delegation (40), judge turns (40, kept small in spirit by
prompting judges to promise few calls), collective turns (40), and a
run-level round ceiling (60, with \texttt{max\_rounds} often set to 20 in
configs). A deterministic stagnation breaker stops the run after
\texttt{stagnation\_rounds} (3) consecutive rounds with no completed task,
after repeated consecutive phase failures, or when a derived token estimate
exceeds the run ceiling; it never consults a model. The fleet analyst owns
a model ladder (weakest first, config-driven) and promotes a role one rung
on repeated failure.

\subsection{Scoring}

A run succeeds only when the evaluator reruns the target toolchain
(\texttt{cargo build} + \texttt{cargo test} for Rust targets,
\texttt{pytest} for Python, syntax checking via \texttt{node --check} for
JavaScript) in the produced workspace and every test passes. Agent
self-reports are recorded but never trusted. This rule exists because the
surveys caught real self-report traps: in the 220926 sweep, three ledger
runs wrote passing per-problem records over empty workspaces; in the
2026-09-19 comparison, the harness compared a string
\texttt{compilation\_status == "pass"} while the schema asked the model for
a boolean, so exit codes carried no signal until the scorer started
re-running the toolchain itself.

\section{Supporting Infrastructure}
\label{sec:infra}

\paragraph{Hashline file tools.}
File tools (\texttt{read}, \texttt{write}, \texttt{edit}) address content
through hashline tags (\texttt{PATH\#TAG}). Edits name the tag they replace;
a stale tag is an invalidation, not a silent overwrite. This gives the
guard a path-shaped argument to check and gives the model a cheap
merge-conflict signal.

\paragraph{Snapcompact.}
When the projected input of the next model call exceeds a per-role token
threshold, the snapcompact hook keeps the most recent messages verbatim,
serializes the discarded history to dense text, renders that text to PNG
frames with a bundled public-domain 8x13 BDF bitmap font, and attaches the
frames as image blocks. The pass is deterministic (same history in, same
frames out) and costs no model call---unlike LLM summarization, which both
spends budget and distorts. Per ADR 0022 the hook attaches to every role,
with lower thresholds for the long-context specialists (translator,
validator) that need it most.

\paragraph{Continuation patch.}
Ollama's chat template for Qwen models rejects a message list that ends
with tool results and carries no trailing user query (HTTP 500, ``no user
query found in messages''). rig 0.42 produces exactly that shape on every
multi-turn tool continuation. A completion-call hook appends a synthetic
user message (``Tool results above. Continue the task from here.'') to the
history; without it each rejected call burns a provider retry of 15--55~s.

\paragraph{Tool sandbox.}
Tool executions run against snapshot-checked workspaces, so a bad write can
be rolled back without touching the source tree.

\section{Experimental Setup}
\label{sec:setup}

\paragraph{Model and hardware.}
All reported runs use \texttt{qwen3.8:27b-mtp-q4\_K\_M} (a 27B reasoning
model with multi-token prediction, 4-bit quantized) served locally by
ollama on one RTX~3090, fully offloaded. A 2B model
(\texttt{minicpm5-2b:q8\_0}) and a 27B fine-tune
(\texttt{Swift-Qwen3.8-27B:dflash2}) appear only in the 220926 sweep.
Typical generation settings: \texttt{num\_ctx} 16384, temperature 0.2,
\texttt{think: true}, one retry.

\paragraph{Problem sets.}
Two families. The \emph{kata} family is the GildedRose-Refactoring-Kata
(6.8~KB of C, translated to Rust) used for the method-comparison protocol.
The \emph{repository} family is the ReCodeAgent dataset: language-pair
directories under \texttt{assets/ReCodeAgent/data/tool\_projects} (c$\to$rust,
python$\to$javascript, go$\to$rust, java$\to$python), with per-project
configs. The autooptimise cycle-1 set spans six pairs from 125 to
$\sim$28k source LOC (Table~\ref{tab:problems}).

\begin{table}[ht]
\centering
\begin{tabular}{llll}
\hline
Problem set & Source & Target & Source LOC \\
\hline
crust/nandc & C & Rust & 125 \\
skel/colorsys & Python & JavaScript & 355 \\
crust/utf8 & C & Rust & 806 \\
crust/libm17 & C & Rust & 7{,}009 \\
oxidizer/stats & Go & Rust & 5{,}625 \\
alphatrans/commons-validator & Java & Python & 28{,}110 \\
\hline
\end{tabular}
\caption{The cycle-1 autooptimise problem sets (measured on source trees;
headers excluded for C). Test commands come from the config: \texttt{cargo
test}, \texttt{pytest}, or \texttt{node --check}.}
\label{tab:problems}
\end{table}

\paragraph{Budgets.}
The MAS configs fix \texttt{max\_rounds} 20, manager turns 20 per round,
worker turns 40 per delegation, judge turns 40, fanout 2, stagnation
breaker at 3 rounds, \texttt{plan\_cap} 4000, \texttt{notes\_cap} 8000.
Baselines used \texttt{max\_output\_tokens} 4096; the H1 candidate raises
it to 8192 (Section~\ref{sec:frontier}).

\paragraph{Protocol.}
One process per cell, strictly sequential, fresh workspace per run, model
offloaded between runs and verified by polling the daemon. Every run leaves
an immutable experiment directory per the ADR 0020 contract:
\texttt{manifest.json} (model, budgets, git revision, parents, hypothesis),
config, \texttt{traces/turns.jsonl} (one line per model call, tool call, and
tool result), the produced \texttt{workspace/}, captured streams, and a
\texttt{result/} aggregate plus per-problem records with a staged failure
taxonomy (\texttt{setup}, \texttt{analyze}, \texttt{plan},
\texttt{translate}, \texttt{validate}, \texttt{evaluate}).

\section{Results}
\label{sec:results}

\subsection{Method comparison on the kata (2026-09-19)}

Three paper-faithful methods on the GildedRose C$\to$Rust kata, 5 passes
each, toolchain-only scoring (Table~\ref{tab:kata}).

\begin{table}[ht]
\centering
\begin{tabular}{lllll}
\hline
Method & Compiled & Full pass & Test pass rate & Wall \\
\hline
Monolith (ReAct + validator) & 5/5 & 1/5 & $0.50 \pm 0.00$ & $97.4 \pm 34.0$ s \\
Ledger~\cite{ledger2026} & 5/5 & 1/5 & $0.60 \pm 0.22$ & $116.3 \pm 21.6$ s \\
ReCodeAgent~\cite{recodeagent2026} & 5/5 & 1/5 & $0.65 \pm 0.22$ & $160.5 \pm 19.1$ s \\
\hline
\end{tabular}
\caption{Kata benchmark, revision 2 (paper-faithful implementations, 5
passes per method). The two multi-agent methods sit above the monolith
(0.65, 0.60 vs.\ 0.50), but the deltas sit inside one SD; the evidence for
scaffolds paying is one full pass in five for each method. An earlier
flattened-implementation generation had the monolith on top (0.75); the v1
and v2 monolith columns are two harness generations, not one method twice.}
\label{tab:kata}
\end{table}

Two readings carry forward. First, on a small task with a 27B model, the
single agent is already at or above the regime where decomposition pays,
and orchestration turns subtract from the fix budget---the papers' own
conclusion (scaffolds pay most where the model is weakest) predicts exactly
this inversion. Second, wall time ranks stably with agent count: more
agents, more seconds.

\subsection{Survey 220926: capability dominates methodology (180 cells)}

Grid: 4 projects (233--1110 C LOC) $\times$ 3 methods $\times$ 3 models
$\times$ 5 reps. Toolchain-verified end-to-end success: 5 of 180 runs
(2.8\%), all on the 233-LOC project, all monolith, all on a 27B model. The
2B model never produced a compiling workspace. Ledger wrote 3 per-problem
records, but every one was a manager self-report over an empty
workspace---caught by post-sweep inspection, and the direct motivation for
toolchain-only scoring. Recode never produced a record. The failure
taxonomy (Table~\ref{tab:failures220926}) put turn-budget exhaustion at
42\% and daemon-rejected tool-call arguments at 13\%.

\begin{table}[ht]
\centering
\begin{tabular}{lrr}
\hline
Failure class & Count & Share \\
\hline
MaxTurns (budget exhausted mid-translation) & 75 & 42\% \\
Other (model-side protocol errors, empty streams) & 70 & 39\% \\
Invalid tool-call arguments (daemon HTTP 500) & 23 & 13\% \\
JSON parse of model output & 4 & 2\% \\
Toolchain-verified success & 5 & 3\% \\
Self-reported ``pass'' over empty workspace & 3 & 2\% \\
\hline
\end{tabular}
\caption{Failure taxonomy over 180 runs (survey 220926). The multi-agent
scaffolds divide the same turn budget across more workers, so they exhaust
it sooner relative to useful work.}
\label{tab:failures220926}
\end{table}

\subsection{Survey 230926: budgets unbind, size binds (105 cells)}

With generous budgets (120 turns, 8192 output tokens, 2-hour timeouts),
MaxTurnsError went from 42\% of runs to zero; mean turns used stayed under
16 per method. On 7 projects (185--3087 C LOC) $\times$ 3 methods
$\times$ 5 reps with qwen3.8-27B, 39 of 105 cells (37\%) verified. Success
concentrated by project size, not method: the 185--324 LOC band converted
at 87--93\% per attempt; the 635--3087 LOC band converted at 0\% (nothing
compiled), with a sharp transition between 324 and 635 LOC. Method
differences were second-order and project-dependent: ledger was the only
method to pass \texttt{chtrie} (5/5), monolith the only one to pass
\texttt{murmurhash\_c} (5/5), recode never won a cell outright but still
passed 4/5 on two projects. Ledger used the fewest turns (mean 5.3) but the
longest walls (mean 630~s); monolith used mean 11.8 turns and 316~s. The
220926 ledger self-report trap did not recur: all 15 ledger scored runs
carried real evaluator steps.

\subsection{The frontier loop: baselines stall at Contract}
\label{sec:frontier}

The autooptimise loop (ADR 0020, after Meta-Harness~\cite{metaharness2026})
runs one immutable experiment directory per candidate, with the proposer
reading raw traces, keeping a Pareto frontier over (pass rate, tests
passed, tests failed), and never editing \texttt{scripts/} or tuning
against the problem set.

\paragraph{Baseline stagnation.}
Every baseline MAS run across six problem sets aborted identically:
3 consecutive stalled rounds, final phase \textsf{Discovery} or
\textsf{Contract}, no translated project. The colorsys baseline
(20260926T213500Z) reached Contract and had its brief rejected
(43 lines against a 40-line cap); the fix delegation then burned its entire
4096-token output budget on a single think-only turn
(\texttt{finish\_reason=Length}, empty text, no tool call), the provider
raised ``produced no answer,'' the round stalled, and the breaker tripped.
The crust-nandc baseline showed the identical signature. The cross-run
invariant: \emph{every} delegation that died did so with output tokens
exactly equal to the 4096 cap and empty visible text---the reasoning block
consumed the whole budget.

\paragraph{H1 diagnosis and fix.}
The mechanism chain: the config enables \texttt{think: true}; the 27B model
emits long reasoning (\texttt{think}) blocks; rig counts them against
\texttt{max\_output\_tokens}; and a turn that delivers no answer and is
truncated (\texttt{finish\_reason=Length}) is a hard completion error, not
a retry with more budget. The harness treats the delegation as failed, the
manager retries the same shape, the model re-thinks just as long, and three
stalls trip the breaker. H1 proposed the minimal mechanism-targeted change:
raise \texttt{max\_output\_tokens} 4096 $\rightarrow$ 8192 (the library
default; the configs had pinned 4096 explicitly). Alternatives considered
and deferred: \texttt{think: false} (kills deliberate reasoning), a
continuation-on-length hook (needs provider internals), and loosening the
validator's line cap (symptom-level, against the leakage rules).

\paragraph{Effect.}
The H1 candidate on crust/nandc (20260927T112000Z, parents = the two
nandc baselines, hypothesis recorded in the manifest) completed 20 rounds,
13 delegations, 12 of 13 tasks done, and reached \textsf{Hardening}---the
first MAS run to leave the Contract phase---and its workspace passed the
toolchain rerun with 2 tests passed, 0 failed. This is the first
translated, toolchain-verified repository migration produced by the MAS
method. It remains incomplete (the run ended at Hardening, not Done), and
the other five H1 candidates in cycle 1 still stalled at Discovery or
Contract; the fix unblocked the mechanism, not every run. H2 (queued for
cycle 2) targets the second observed mechanism: a first-call request shape
with an empty prompt field that ollama's server-side chat template rejects
with HTTP 500, which the continuation hook does not yet cover; 23 of 28
model calls in the base-crust-libm17 trace carried the empty prompt.

\paragraph{Frontier discipline.}
The frontier keeps non-dominated candidates rather than a champion: the
H1-nandc candidate dominates the nandc baselines on pass rate and tests
passed; the other five sets have no scored candidate yet, so their
baselines stay on the frontier by default. Lineage lives in the manifests
(\texttt{parents}), and every hypothesis is one causal sentence recorded
before the run.

\section{Discussion}
\label{sec:discussion}

\paragraph{Failure modes.}
Three failure modes dominate everything we have measured.
(1)~\emph{Think-truncation}: reasoning blocks consume the output budget;
the turn dies with \texttt{finish\_reason=Length} and no visible answer;
retries reproduce it deterministically at temperature 0.2. Fixed by budget,
not by prompt surgery. (2)~\emph{Provider shape rejections}: ollama's
server-side chat template 500s on message lists without a trailing user
query, both on tool continuations (patched by the hook) and on a first-call
shape (H2, open). (3)~\emph{Stagnation}: without a breaker, the manager
retries the same delegation shape indefinitely; the 3-round breaker
converts an unbounded loop into a diagnosable artifact. A fourth mode,
malformed tool-call JSON (13\% of survey 220926), largely disappeared from
the MAS runs after decisions moved from tool calls to parsed text.

\paragraph{Limitations.}
The verified-success frontier for a 27B model on C$\to$Rust sits near
300--350 source LOC for single-shot methods; MAS pushed one 125-LOC
repository through Contract to Hardening but no run has yet landed
\textsf{Done} end-to-end. Fan-out is turn-level interleaving, not
parallelism, while the fleet has one model slot. Autooptimise has completed
one proposal--evaluate cycle; the leakage rule (holdout problem sets must
differ from search sets) blocks any generalization claim until a second
family is scored. The alphatrans set (28k LOC Java) exceeds what a 16k
context can even inventory; it stands as a stress upper bound, not a
current target.

\paragraph{Threats to validity.}
Sample sizes are small ($n{=}5$ per cell in the kata benchmark; single
completions per problem set in the frontier loop), so no significance
claims are made; deltas inside one SD are reported as such. One model and
one hardware configuration dominate the data; the 220926 model-axis
evidence is three models, one of which never worked. Infra failures
contaminate raw counts (57 of 109 on-disk 230926 attempts were preflight
aborts against a dead daemon) and are excluded from per-method rates, which
the experiment contract now prevents by gating on daemon liveness.
Finally, toolchain-only scoring can still pass a workspace whose tests are
weaker than the source suite; the tester's per-module generated-test cap
(4) bounds but does not eliminate that gap.

\section{Conclusion}

ARCMiS shows that repository-level code migration with small local models
is tractable only when the harness carries most of the discipline: evidence-
gated phases, a deterministic guard on every tool call, character-capped
shared state, toolchain-only scoring, and a stagnation breaker that turns
silent loops into artifacts. The measured record is deliberately honest:
the simplest method won the first two surveys, capability dominated
methodology at the 2B/27B boundary, project size between 324 and 635 LOC
marked a sharp capability wall, and the MAS method's first verified
migration came from a one-line output-budget fix that the optimization loop
found by reading raw traces---not from a cleverer prompt. The infrastructure
(the experiment contract, the frontier discipline, the failure taxonomy) is
the durable contribution; it makes every next hypothesis cheap to test and
every claimed success impossible to fake.

\bibliographystyle{plainnat}
\bibliography{references}

\end{document}
